\subsection{Benchmark data}
Reference coding regions were taken from the RefSeq \citep{OLeary2016} transcript records (\texttt{rna.gbff}) of eight
eukaryotes (Table~\ref{tab:orgs}):
\begin{itemize}
  \item the four organisms of the original 2018 evaluation;
  \item human, restricted to MANE Select transcripts \citep{Morales2022};
  \item \emph{Tetrahymena thermophila}, which uses the ciliate genetic code: TAA and TAG encode
        glutamine, and TGA is the only stop;
  \item two compositional extremes: \emph{Plasmodium falciparum} (\todo{x}\% GC) and
        \emph{Chlamydomonas reinhardtii} (\todo{y}\% GC).
\end{itemize}
Curated (NM\_) records were used where available; model (XM\_) records were used for the last three.

A transcript was eligible if its single annotated CDS:
\begin{itemize}
  \item was complete and contiguous;
  \item began with ATG and ended with a stop of the record's translation table;
  \item had no internal stop, frameshift, translational exception or ribosomal slippage;
  \item translated to the annotated protein.
\end{itemize}
For each gene the eligible transcript with the longest CDS was kept. UTRs were trimmed to at most
250\,nt, as in the 2018 evaluation.

Each truth set was used to derive four inputs (all generated with fixed seeds):
\begin{enumerate}
  \item \emph{complete}: the transcripts as given;
  \item \emph{fragments}: each transcript is cut once within its CDS at a fraction drawn from
        $\mathcal N(0.5,0.1)$. Cuts 3$'$ of the CDS midpoint keep the 3$'$ piece, which retains the
        stop but not the start; other cuts keep the 5$'$ piece, which retains the start but not the
        stop;
  \item \emph{unstranded} versions of (1) and (2), in which a random half of the sequences are
        reverse-complemented;
  \item each of the above with one random negative sequence per transcript added. A negative has
        the transcript's length and G+C content.
\end{enumerate}

\subsection{Compared methods}
The following were run with default parameters unless stated:
\begin{itemize}
  \item \PyTD, with and without start refinement;
  \item Perl \TD\ (current release, v6.0.0 and v5.7.1);
  \item Prodigal 2.6.3 \citep{Hyatt2010}, in single mode with a full motif scan (\texttt{-n}) and
        the organism's translation table;
  \item CodAn 1.1 \citep{Nachtigall2021}, with the pre-trained model of the closest clade (full or
        partial model to match the input);
  \item \todo{GeneMarkS-T \citep{Tang2015}};
  \item two baselines: the longest ORF per transcript, and all ORFs of at least 100 codons.
\end{itemize}
Inputs labelled strand-specific were analysed on the forward strand only (\TD\ \texttt{-S};
Prodigal forward-strand calls; CodAn \texttt{--strand plus}). Unstranded inputs were analysed on both
strands.

\subsection{Accuracy measures}
Following the 2018 evaluation, a prediction matched a reference CDS on the same transcript and strand
when its 3$'$ end was within 2\,nt of the reference 3$'$ end (stop-level match). It was an exact match
when its 5$'$ end was also within 2\,nt. Each reference CDS could be matched at most once.

We report:
\begin{itemize}
  \item sensitivity (matched references / references);
  \item precision (matched predictions / all predictions, including predictions on random
        negatives);
  \item their harmonic mean (F1).
\end{itemize}
\todo{Runtime measured on \ldots}
